% from ega4-2-fr.tex
\oldpage[IV-2]{143}

\subsection{\texorpdfstring{Platitude et propriété $(R_k)$}{Platitude et propriété (R\string_k)}}

\begin{proposition}[6.5.1]
\EGAReaderTarget{ega4.c04.prop.6-5-1}{IV.2.s6.5.1-fr}
Soient $A$, $B$ deux anneaux locaux noethériens, $k$ le corps résiduel
de $A$, $\varphi:A\to B$ un homomorphisme local, pour lequel $B$ est un
$A$-module plat. Alors:
\begin{enumerate}
\item[(i)] Si $B$ est régulier, $A$ est régulier.
\item[(ii)] Si $A$ et $B\otimes_Ak$ sont réguliers, $B$ est régulier.
\end{enumerate}
Cette proposition est donnée pour mémoire, ayant déjà été démontrée dans
(0, 17.3.3).
\end{proposition}

\begin{corollary}[6.5.2]
\label{IV.2.s6.5.2-fr}
Soient $X$, $Y$ deux préschémas localement noethériens, $f:X\to Y$
un morphisme plat.
\begin{enumerate}
\item[(i)] Si $X$ est régulier en un point $x$, $Y$ est régulier au point
$f(x)$.
\item[(ii)] Si pour un $y\in f(X)$, $Y$ est régulier au point $y$ et si
$f^{-1}(y)$ est un préschéma régulier en un point $x$, alors $X$ est
régulier au point $x$.
\end{enumerate}
\end{corollary}

\begin{proposition}[6.5.3]
\label{IV.2.s6.5.3-fr}
Soient $X$, $Y$ deux préschémas localement noethériens, $f:X\to Y$
un morphisme plat.
\begin{enumerate}
\item[(i)] Si $X$ possède la propriété $(R_k)$ en un point $x$
(resp. si $X$ possède la propriété $(R_k)$), $Y$ possède la propriété
$(R_k)$ au point $f(x)$ (resp. $Y$ possède la propriété $(R_k)$ en tous
les points de $f(X)$).
\item[(ii)] Supposons que pour tout $y\in f(X)$, le préschéma
$f^{-1}(y)$ possède la propriété $(R_k)$; alors, si pour un point
$y\in f(X)$, $Y$ possède la propriété $(R_k)$ au point $y$, $X$ possède
la propriété $(R_k)$ en tout point de $f^{-1}(y)$.
\end{enumerate}
\end{proposition}

\begin{proof}
\medskip\noindent
\begin{enumerate}
\item[(i)] Posons $y=f(x)$ et soit $y'$ une générisation de $y$; comme
dans la démonstration de (6.4.1), il y a un point générique $x'$ d'une
composante irréductible de $f^{-1}(y')$ qui est une générisation de $x$;
donc on a
\[
\dim(\mathcal O_{x'}\otimes_{\mathcal O_{y'}}k(y'))=0
\]
et, en vertu de l'hypothèse et de (6.1.2),
\[
\dim(\mathcal O_{x'})=\dim(\mathcal O_{y'}).
\]
L'hypothèse entraîne, soit que
$\dim(\mathcal O_{x'})>k$, auquel cas
$\dim(\mathcal O_{y'})>k$, soit que $\mathcal O_{x'}$ est un anneau
régulier, et alors $\mathcal O_{y'}$ est un anneau régulier en vertu de
(6.5.1).

\item[(ii)] Comme, pour toute générisation $x'$ de $x\in f^{-1}(y)$,
$y'=f(x')$ est une générisation de $y$, on peut se borner à vérifier que
si $\dim(\mathcal O_x)\leqslant k$, alors $\mathcal O_x$ est un anneau
régulier. Or, on a, en vertu de (6.1.2),
\[
\dim(\mathcal O_x)=\dim(\mathcal O_y)
+\dim(\mathcal O_x\otimes_{\mathcal O_y}k(y)).
\]
Donc si $\dim(\mathcal O_x)\leqslant k$, on a a fortiori
$\dim(\mathcal O_y)\leqslant k$ et
$\dim(\mathcal O_x\otimes_{\mathcal O_y}k(y))\leqslant k$, et l'hypothèse
entraîne que $\mathcal O_y$ et
$\mathcal O_x\otimes_{\mathcal O_y}k(y)$ sont des anneaux réguliers. On
conclut alors de (6.5.1) que $\mathcal O_x$ est un anneau régulier.
\end{enumerate}
\end{proof}

\begin{corollary}[6.5.4]
\label{IV.2.s6.5.4-fr}
Soient $X$, $Y$ deux préschémas localement noethériens, $f:X\to Y$
un morphisme plat.
\begin{enumerate}
\item[(i)] Si $X$ est normal en un point $x$, $Y$ est normal au point
$f(x)$.
\item[(ii)] Si, pour tout $y\in f(X)$, $f^{-1}(y)$ est un préschéma
normal et si $Y$ est normal en un point $y\in f(X)$, alors $X$ est un
préschéma normal en tout point de $f^{-1}(y)$.
\end{enumerate}
\end{corollary}

Cela résulte aussitôt du critère de normalité de Serre (5.8.6) et de
(6.4.1) appliqué pour $k=2$ et de (6.5.3) appliqué pour $k=1$.

